\documentclass[11pt]{article}

\usepackage{amssymb,amscd,amsthm, verbatim,amsmath,color,fancyhdr, mathrsfs}
\usepackage{graphicx}
\usepackage{turnstile}
\usepackage[hidelinks]{hyperref}

\def\tequiv{\ensuremath\sdststile{}{}}

\usepackage[letterpaper, left=2cm, right=2cm, top=2.5cm,
bottom=2.5cm,dvips]{geometry}


\newtheorem{thm}{Theorem}[section]
\newtheorem{prop}[thm]{Proposition}
\newtheorem{lem}[thm]{Lemma}
\newtheorem{rem}[thm]{Remark}
\newtheorem{cor}[thm]{Corollary}
\newtheorem{exr}[thm]{Exercise}
\newtheorem{ex}[thm]{Example}
\newtheorem{defn}[thm]{Definition}

% \underset{x}{\operatorname{argmax}} 
%
% \DeclareMathOperator*{\argmin}{arg\,min}

\newcommand{\R}{\mathbb{R}}
\newcommand{\Z}{\mathbb{Z}}
\newcommand{\N}{\mathbb{N}}

\newcommand{\aA}{\mathcal{A}}
\newcommand{\aB}{\mathcal{B}}
\newcommand{\aM}{\mathcal{M}}
\newcommand{\aL}{\mathcal{L}}

\title{Equations de Hamilton Jacobi}

\author{Victor Schmit}

\date{\today}

\begin{document}
\maketitle

\tableofcontents
\newpage

\section{Transformée de Legendre}

\defn[Transformée de Legendre]
Soit $f : \R^d  \to \R$ une fonction convexe, superlinéaire. On définit la transformée
de Legendre $\aL$ de $f$ par
\[
	\aL(f) (x) = \max_{y \in \R^d} \langle x \cdot y\rangle - f(y).
\]

\rem Comme $\R^d$ est de dimension finie, on peut montrer que $f$ est continue.

\rem On verra que la transformée de Legendre fait apparaître des notions de dualité :
une notation utilisée pour la transformée de $f$ est $f^* = \aL(f)$.

\rem La tansformée de Legendre est bien définie. En effet, par superlinéairité,
\[
	\lim_{\|y\| \to \infty} \frac{f(y)}{\|y\|} = + \infty.
\]
Donc, pour des raisons de continuité et de coercivité,
\[
	y \mapsto - \langle x \cdot y \rangle + f(y) \quad \text{admet un minimum.}
\]
La transformée de Legendre est donc bien définie.

\prop Quelques propriétés sur la transformée de Legendre.
\begin{enumerate}
	\item $\aL(f)$ est convexe.

	\item $\aL(f)$ est superlinéaire.

	\item $\aL$ est une involution.
\end{enumerate}

\begin{proof}
	Soit $f : \R^d \to \R$ une fonction convexe superlinaire.

	\medskip

	\emph{(1)} Montrons que $f^*$ est convexe. Soit $x,y,u \in \R^d$ et $\lambda  \in [0,1]$ et $z  = x + \lambda \cdot (y-x)$.
	Alors,
	\[
		f^*(z) \le \langle u \cdot z\rangle - f(u) = (1-\lambda) \cdot (\langle u \cdot x \rangle - f(u)) +
		\lambda \cdot (\langle u \cdot y \rangle - f(u)).
	\]
	Il est donc clair que
	\[
		f^*(z) \le (1-\lambda) \cdot f^*(x) + \lambda \cdot f^*(y).
	\]
	On en déduit que $f^*$ est convexe.

	\medskip

	\emph{(2)} Montrons que $f^*$ est superlinéaire. Soit $x \in \R^d$ et $A > 0$. On pose
	$y = A \cdot \frac{x}{\|x\|}$. Dans ce cas, on a
	\[
		f^*(x) \ge A \|x\| - f(A \cdot \frac{x}{\|x\|}).
	\]
	Donc par continuité de $f$ et compacité des boules fermés,
	\[
		\frac{f^*(x)}{\|x\|} \ge A  - \frac{1}{\|x\|}\max_{u\in B(0,A)}f(u).
	\]
	On en déduit que
	\[
		\liminf_{\|x\| \to +\infty} \frac{f^*(x)}{\|x\|} \ge A.
	\]
	D'où, en faisant tendre $A$ vers l'infini,
	\[
		\liminf_{\|x\| \to +\infty} \frac{f^*(x)}{\|x\|} \ge +\infty.
	\]
	On en déduit que la limite en l'infini existe et que
	\[
		\lim_{\|x\| \to +\infty} \frac{f^*(x)}{\|x\|} = +\infty.
	\]
	On en conclut que $f^*$ est superlinéire.

	\medskip

	\emph{(3)} Montrons que $\aL$ est une involution. On vient de montrer que $\aL$ envoie son espace de définition
	dans son espace de définiton. Il est donc tout à fait naturel de composer successivement une telle transformée. On montre ce résultat par double inégalités.

	\emph{Etape 1 :} Montrons que $\aL^2 f \le f$. Soit $x \in \R^d$. On fixe $y$ tel que
	$\aL^2f(x) = \langle x \cdot y \rangle -\aL f (y)$. Dans ce cas, on prend un $z \in \R^d$ libre. On a
	\[
		\aL^2 f (x) = \langle x \cdot y \rangle -\aL f (y) \le \langle x - z \cdot y \rangle  +  f (z)
	\]
	En prenant $z = x$, on a bien $\aL^2 f (x) \le f(x)$.

	\medskip

	\emph{Etape 2 :} Montrons que $\aL^2 f \ge f$. Soit $x \in \R^d$. Cette fois, on laisse $y \in \R^d$ libre. On a
	\[
		\aL^2 f (x) \ge \langle x \cdot y \rangle -\aL f (y).
	\]
	On peut donc choisir $z \in \R^d$ de telle sorte que
	\[
		\aL f (y) = \langle z \cdot y \rangle - f (z).
	\]
	On en déduit que
	\[
		\aL^2 f (x) \ge \langle x-z \cdot y \rangle + f (z).
	\]
	Pour la suite, il faut choisir $y$ soigneusement. Comme $f$ est convexe, il existe $y \in \R^d$ tel que pour tout $u \in \R^d$ on aie
	\[
		\langle u-x \cdot y \rangle \le f(u) - f(x)
	\]
	Dans ce cas, il est clair que $\aL f (y) = \langle x \cdot y \rangle - f(x)$. Donc, en posant $z = x$ on a
	\[
		\aL^2 f (x) \ge f(x).
	\]
	Cela permet de conclure.
\end{proof}

\rem Cette transformation est très intéressante car elle permet de faire le lien entre le Lagrangien et le Hamiltonien.
Pour cela, il suffit de définir $L$ comme la transformée de Legendre de du hamiltonien $H$.

\ex Si on prend  $f(x) = \frac{\|x\|^2}{2}$, alors $f^* = f$
\ex Si $f$ est différentiable et que $f^*(z) = \langle z \cdot x \rangle - f(z)$, alors
$df(z) = x$.

\section{Formule de Hopf Lax}

On propose de résoudre l'équation
\begin{equation}
	\begin{cases}
		\partial_t u + H(\nabla_x u) = 0, \\
		u(0,x) = u_0(x)
	\end{cases}
	\label{eq:ham}
\end{equation}
avec $H$ convexe, superlinéaire et u globalement lipschitzienne.

\medskip

On note $L$ la transformée de Legendre du Hamiltonien $H$ que l'on nommera le Lagrangien. On revient à la formulation variationnelle :

\begin{equation}
	u(t,x) = \inf_{\substack{\omega \in \mathcal{A}_{t,x,y} \\
			y \in \mathbb{R}^d}}
	\left(
	\int_{0}^t L(\dot{\omega}(s))\, ds + u_0(y)
	\right)
	\label{eq:lag}
\end{equation}
Si $u$ est une solution de l'équation eq.~\ref{eq:lag}, alors $u$ est une solution de l'équation eq.~\ref{eq:ham}.

\thm[Formule de Hopf Lax] Si on suppose que $L$ est convexe et superlinéaire et que $u_0$ et globalement lipschitzienne, alors si $u$ satisfait l'équation eq.\ref{eq:lag}, on a pour tout $t > 0$ et $x \in \R^d$

\begin{equation}
	u(t,x) = \min_{y \in \R^d} t\cdot L\left(\frac{x-y}{t}\right) + u_0(y)
	\label{eq:hopf}
\end{equation}

\begin{proof}
	Soit $x,y \in \R^d$

	\medskip

	\emph{Étape 1 : } Montrons que le minimum existe. Pour cela on pose

	\[
		h(y) = t \cdot L \left(\frac{x-y}{t}\right) + u_0(y)
	\]

	Il est clair que $h$ est continue. Il reste donc à montrer que $h$ est coercive pour prouver l'existence d'un minimum. Or, si on note $C$ la cosntante de lipschitz de $u_0$, on a
	\[
		h(y) \ge \|y\| \cdot \left (
		\frac{t}{\|y\|} \cdot L\left(\frac{x-y}{t}\right) -  C
		\right).
	\]
	Or, comme  $\|y\| \sim_{y \to +\infty} \frac{\|x-y\|}{t}$ et comme $L$ est superlinéaire d'après les propriétés de la transformée de Legendre, il devient clair que
	\[
		\frac{h(y)}{\|y\|} \xrightarrow[y \to +\infty]{} +\infty.
	\]
	Donc, $h$ est superlinéaire et en particulier est coercive.

	\medskip

	\emph{Étape 2 : } Soit $\omega \in \aA_{t,x,y}$ un chemin de classe $C^1$ reliant $y$ et $x$ sur $[0,t]$. D'après les propriétés de la transformée de Legendre,
	le Lagrangienn $L$ est convexe. On peut donc utiliser l'inégalité de Jensen pour les fonctions convexes. On a
	\[
		\displaystyle\int_{[0,t]}^{} L(\dot{\omega}(s)) ds \ge
		t \cdot L\left(\frac{1}{t}\displaystyle\int_{[0,t]}^{}   \dot{\omega}(s) ds\right)
		= t \cdot L\left(\frac{x-y}{t}\right)
	\]
	Donc,
	\[
		\displaystyle\int_{[0,t]}^{} L(\dot{\omega}(s)) ds + u_0(y) \ge
		t \cdot L\left(\frac{x-y}{t}\right) + u_0(y)
	\]
	En passant, aux bornes inférieures, on a donc
	\[
		u(t,x) \ge
		\min_{y \in \R^d} \left( t \cdot L\left(\frac{x-y}{t}\right) + u_0(y) \right)
	\]

	\medskip

	\emph{Étape 3 : } Pour montrer l'autre inégalité, on choisit $y$ qui minimise $h$. Dans ce cas, on pose l'impulsion $v = \frac{x-y}{t}$. Dans ce cas, on peut définir
	$\omega : [0,t] \to \R^d$ par
	\[
		\omega(s) = s \cdot v + u_0(y).
	\]
	Il est clair que $\omega \in \aA_{t,x,y}$. De plus,
	\[
		\displaystyle\int_{[0,t]}^{} L(\dot{\omega}(s)) ds + u_0(y) =
		\displaystyle\int_{[0,t]}^{} L(v) ds + u_0(y) =
		t \cdot L\left(v\right) + u_0(y).
	\]
	Comme on a choisi $y \in \R^d$ pour minimiser $h$, on a donc
	\[
		\displaystyle\int_{[0,t]}^{} L(\dot{\omega}(s)) ds + u_0(y) =
		\min_{y \in \R^d} \left( t \cdot L\left(\frac{x-y}{t}\right) + u_0(y) \right).
	\]
	On en conclut que
	\[
		\min_{y \in \R^d} \left( t \cdot L\left(\frac{x-y}{t}\right) + u_0(y) \right)
		\ge
		u(t,x).
	\]
\end{proof}

\prop On pose la fonction $u : \R^*_+ \times  \R^d \to \R$ définie par la fomule de Hopf eq.~\ref{eq:hopf}. Dans ce cas,
\[
	u(t,x) = \min_{z \in \R^d}\left(
	(t-s) \cdot L \left(\frac{x-z}{t-s}\right) + u(s,z)
	\right)
\]
De plus, $u$ est globalement lipschitzienne.

\begin{proof}
	Soit $t > s > 0$ et $x,y,z \in \R^d$.

	\medskip

	Avant de prouver un à un les différents points, on met en évidence une inégalité qui sera utile. On pose les vecteurs vitesse suivants :
	\[
		u = \frac{x-z}{t-s}, \quad v = \frac{z-y}{s}, \quad \omega = \frac{x-y}{t}.
	\]
	Dans ce cas, il est clair que
	\[
		x-y = (t-s) \cdot u + s \cdot v.
	\]
	Donc,
	\[
		\omega = \frac{t-s}{t} \cdot u + \frac{s}{t}\cdot v.
	\]
	On en déduit que

	\[
		t \cdot L(\omega) + u_0(y) \le
		(t-s) \cdot L(u) + s \cdot L(v) + u_0(y).
	\]

	\medskip

	\emph{Étape 1 : } On fixe $x \in \R^d$ et $t>s>0$.
	Dans ce cas,
	\[
		u(t,x) \le t \cdot L(\omega) + u_0(y) \le
		(t-s) \cdot L(u) + s \cdot L(v) + u_0(y).
	\]
	On peut donc ensuite chosir un $y \in \R^d$ tel que
	\[
		u(t,x) \le
		(t-s) \cdot L(u) + u(s,z)
	\]
	Cela étant pour un $z \in \R^d$ quelconque, on a bien
	\[
		u(t,x) \le
		\min_{z\in \R^d} \left((t-s) \cdot L\left(\frac{x-z}{t-s}\right) + u(s,z) \right)
	\]

	\medskip

	\emph{Étape 2 : } Pour montrer l'autre inégalité, on remarque que l'on peut
	choisir un $z \in \R^d$ tel que tous les vecteurs vitesse $v,u,\omega$ soient égaux.
	Pour cela, on peut prendre $z = \frac{s}{t} x + \frac{t-s}{t} y$. On en déduit que,
	\[
		\min_{z\in \R^d} \left((t-s) \cdot L\left(\frac{x-z}{t-s}\right) + u(s,z) \right)\le
		(t-s) \cdot L(u) + s \cdot L(v) + u_0(y)=
		t \cdot L(\omega) + u_0(y).
	\]
	En prenant $y \in \R^d$ minimisant le côté droit, on a
	\[
		\min_{z\in \R^d} \left((t-s) \cdot L\left(\frac{x-z}{t-s}\right) + u(s,z) \right)\le
		u(t,x).
	\]

	\emph{Étape 3 : } Montrons que $u$ est globalement lipschitzienne.


\end{proof}

\end{document}
